% Source comparison and re-derivation for Paper I's cusp-origin calculation.
\begin{lemma}[The actual coherent cusp origin]
\label{mapsaudit:coherent-cusp-origin}
In the fixed exponential cusp chart let $q=e^{2\pi i s}$,
$\Pi_c\widehat\gamma=(6\mu,-s+C_{c,21})^t$,
$C_{c,21}=c+C_{0,21}$, and
$\ell=(2\pi i)^{-1}\log(F/G)=-s/12+\kappa_0(q)$.
Put $D=C_{0,21}-12\kappa_0$ and
\[
 \sigma(c)=\left(-\frac{c\mu}{2},-\frac{c^2}{24}-\frac{cD}{12}\right).
\]
For the coherent even-shift representative
$b_{h,c}=-h\Pi_c\widehat\gamma/12+(h\ell-h^2/24)e$,
its actual holomorphic cusp logarithm is
\[
 \overline b_{h,c}
 =\left(-\frac{h\mu}{2},-\frac{h(c+D)}{12}-\frac{h^2}{24}\right)
 =\sigma(c+h)-\sigma(c).
\]
Thus translating the logarithmic cusp coordinates by $-\sigma(c)$
trivializes this coherent action. This assertion is for this explicit
representative and this fixed exponential chart.
\end{lemma}
\begin{proof}
Substitution of the two displayed ordinary-to-cusp formulas cancels
$hs/12$ against $-hs/12$ in the second component. It gives
$-hC_{c,21}/12+h\kappa_0-h^2/24=-h(c+D)/12-h^2/24$.
The first component is $-h\mu/2$. Direct subtraction of the two
quadratic polynomials $\sigma(c+h)$ and $\sigma(c)$ gives precisely
this vector. Its entries are holomorphic on the full cusp disc.
Exponentiation therefore acts in every toric chart by multiplication
by nonzero holomorphic character values; this action and its inverse
extend to every toric stratum. For integral $h$, source and target
lattice multipliers differ by $(1,e^{2\pi i h\lambda_1})=(1,1)$.
The unit action commutes with the integral toric shears and therefore
descends. On the punctured cusp it equals the given ordinary translation,
by the same exponential-chart identity. This proves the origin comparison
on the actual family, beyond the formal cocycle identity.
\end{proof}
